\section{Ensembling: when does combining guards recover transfer?}
\label{app:ensembling}

Ensembling is a classic way to make a classifier generalize, so it is natural to ask whether it repairs
the transfer loss that fine-tuning induces (\Cref{sec:actI}). We test this \emph{retrospectively} on the
adaptation panel's committed per-row margins (\Cref{sec:adaptation}): \EnsNCk{} checkpoints across \EnsNFam{}
families (excluding the degenerate Llama-Guard null cell), each scored under \{unmodified base, SFT,
KL-SFT\} $\times$ five seeds on the \emph{same} rows. The one exception is the cross-model committee of
\Cref{tab:ensembling-committee}, which is reported over the four general checkpoints and over all ten
including the null cell, because a rank-averaged committee is defined by its membership and dropping a
member would change the object rather than clean it. Every number is equal-family macro-AP on the raw
logit margin, anchored to each checkpoint's \emph{own} unmodified base --- identical estimand to the main
text. Within a checkpoint the base and its adapters share one output head, so raw margins are scale-comparable
in \emph{units}; they are not comparable in \emph{dispersion}, and averaging them raw is therefore not
an equal-weight average of judgments. On this panel the base's margins are $2$--$4\times$ wider than its
adapters' (e.g.\ Qwen3-4B $\mathrm{sd}=16.3$ vs.\ $4.2$), so a raw $\tfrac12/\tfrac12$ average is
effectively base-weighted --- one reason the main text's operator calibrates first (\Cref{eq:comp}).
Across checkpoints margins are not comparable at all, so the cross-model committee (below)
rank-normalizes first. This appendix is retrospective on the inspected panel, not a
preregistered claim.

\paragraph{Ensembling a fine-tune with itself is only denoising.} Averaging the five SFT seeds into one
guard raises transfer by a bootstrap-significant but small \EnsSeedGainTrans{} (one-sided LCB
\EnsSeedGainTransLCB{}) over a single SFT run --- yet transfer still lands \emph{below} the base
(\EnsSeedSFTdTrans{} vs.\ base; \Cref{tab:ensembling}). Seed-ensembling KL-SFT behaves the same
(\EnsSeedKLdTrans{}). The reason is structural: every seed is trained by the same recipe on the same data,
so all members share the \emph{same} specialization bias. Averaging them cancels seed \emph{variance}, not
the shared \emph{bias} toward the represented sources --- so it polishes a specialized guard without
un-specializing it.

\paragraph{Ensembling across the specialization axis recovers transfer.} The one ensemble that lifts
transfer above the base is the base $\oplus$ adapter combination --- the same base-plus-adapter idea as
Act~II, but \textbf{not} the promoted operator of \Cref{eq:comp}: this row averages \emph{raw margins}
(the non-promotable ablation of \Cref{sec:actIII-ablations}) and its adapter member is a \emph{five-seed
mean}, so it costs six forward passes rather than two. On this panel it raises transfer by
\EnsBaseSFTvsBaseTrans{} over the base
(one-sided LCB \EnsBaseSFTvsBaseTransLCB{} $>0$), recovering \EnsBaseSFTvsSFTTrans{} relative to the
\emph{five-seed SFT ensemble} (not the single-run SFT row printed in the same table),
at a represented cost of \EnsBaseSFTvsSFTRep{} (95\% CI $[\EnsBaseSFTvsSFTRepLCB, \EnsBaseSFTvsSFTRepUCB]$).
It is a recovery, not a Pareto win: for the strongest base the composed guard still gives back a little
transfer (the Act~II caveat, \Cref{sec:actIII}). \textbf{Cost-matched, the effect is about half as large.}
Recomputing the honest 2-pass analogues on this same \EnsNCk{}-checkpoint panel, as an
equal-\emph{checkpoint} mean (the tabulated row is an equal-\emph{family} mean, so the two weightings
do not coincide and the control below is read as a magnitude check, not a reproduction):
base $\oplus$ \emph{one} adapter on raw margins gains
$+0.012$, and base $\oplus$ one adapter \emph{calibrated} --- exactly \Cref{eq:comp} --- gains $+0.008$,
against the $+0.014$ of the six-pass row as tabulated. So the qualitative conclusion (only crossing the specialization
axis lifts transfer above base) survives at the promoted operator, but the magnitude does not: the
cost-matched gain sits at the one-sided LCB the six-pass row reports, and no main-text claim rests on this
appendix. A generated, bootstrapped row for the 2-pass calibrated operator is a stated gap
(\Cref{sec:roadmap}). \Cref{fig:ensembling-plane} makes the distinction visual:
the seed-ensembling arrows stay inside the ``transfer below base'' band, while the arrow to base $\oplus$
SFT crosses into ``transfer recovered.'' Adding KL-SFT as a third member (base $\oplus$ SFT $\oplus$
KL-SFT) lands essentially on top of base $\oplus$ SFT --- no further gain. The active ingredient is
\emph{diversity across the specialization axis}: the un-tuned base makes off-distribution errors that are
decorrelated from the tuned member's, so averaging cancels them; two specialized members do not. This is
measured directly, not merely inferred from the AP midpoint: on transfer rows the base's per-row errors
correlate only \EnsErrCorrBaseSFT{} with its own fine-tune's, versus \EnsErrCorrSFTSFT{} between two
fine-tune seeds --- the base is far more complementary to the fine-tune than one fine-tune run is to
another, which is exactly why base $\oplus$ SFT recovers transfer while SFT $\oplus$ SFT (below) does not.

\input{generated/tab_ensembling_gen}

\begin{figure}[t]\centering
\includegraphics[width=0.95\linewidth]{fig_ensembling_plane.pdf}
\caption{\textbf{The ensembling plane} (equal-family mean, \EnsNCk{}-checkpoint panel). Change in
represented (x) and held-out transfer (y) macro-AP vs.\ each checkpoint's own base. Ensembling a fine-tune
with itself (orange, single~$\to$~5 seeds) stays in the ``transfer below base'' band; only crossing the
specialization axis by adding the un-tuned base (green arrow, to base~$\oplus$~SFT) lifts transfer above
the base. \mbox{base~$\oplus$~SFT~$\oplus$~KL-SFT} coincides with base~$\oplus$~SFT (adding a second
specialized member buys nothing).}
\label{fig:ensembling-plane}
\end{figure}

\paragraph{A committee of fine-tuned guards does not generalize better.} The strongest form of the
technique --- a cross-model \emph{committee} that rank-averages different guards on the same row --- makes
the point sharpest (\Cref{tab:ensembling-committee}). A committee of the SFT guards is \emph{worse} than
the single best guard on transfer (\EnsCommSFTvsBestGen{} over the four general checkpoints,
\EnsCommSFTvsBestAll{} over all ten): SFT drives every member toward the same represented sources, so their
held-out errors are correlated and averaging cannot cancel them --- it only dilutes the best member. The
committee of \emph{KL-SFT} guards merely breaks even (\EnsCommKLvsBestGen{} / \EnsCommKLvsBestAll{}),
because KL regularization leaves the members less specialized and therefore more diverse --- but that is
the KL penalty doing the work, not the ensemble.

\input{generated/tab_ensembling_committee}

\begin{takeaway}{Does ensembling help here?}
Only when it injects the diversity fine-tuning removed. Keeping the \emph{un-tuned base} in the ensemble
(Act~II composition) recovers transfer above base; ensembling fine-tunes \emph{with each other} --- across
seeds or as a committee of tuned guards --- does not, because specialization correlates their
off-distribution errors. Ensembling is a variance tool, not a bias fix: it costs $N$ inference passes,
does not dissolve the represented/transfer tradeoff, and every composed candidate must still be
recalibrated and gated on both splits like any other (\Cref{sec:practitioner}).
\end{takeaway}
